\documentclass[12pt,a4paper]{article}

\usepackage[utf8]{inputenc}
\usepackage[T1]{fontenc}
\usepackage[french]{babel}
\usepackage{geometry}
\usepackage{graphicx}
\usepackage{hyperref}
\usepackage{listings}
\usepackage{xcolor}
\usepackage{booktabs}
\usepackage{tabularx}
\usepackage{fancyhdr}
\usepackage{titlesec}
\usepackage{enumitem}
\usepackage{amsmath}
\usepackage{float}

\geometry{
  top=2.5cm,
  bottom=2.5cm,
  left=2.5cm,
  right=2.5cm
}

% ── Couleurs ──────────────────────────────────────────────────────────────────
\definecolor{codebg}{RGB}{30,30,30}
\definecolor{codetext}{RGB}{220,220,220}
\definecolor{keyword}{RGB}{86,156,214}
\definecolor{string}{RGB}{206,145,120}
\definecolor{comment}{RGB}{106,153,85}
\definecolor{primary}{RGB}{0,120,212}
\definecolor{darkgray}{RGB}{50,50,50}
% Terminal Ubuntu
\definecolor{termbg}{RGB}{48,10,36}
\definecolor{termfg}{RGB}{211,215,207}
\definecolor{termgreen}{RGB}{78,154,6}
\definecolor{termcyan}{RGB}{52,177,186}
\definecolor{termtitle}{RGB}{30,5,22}

% ── Style listings ────────────────────────────────────────────────────────────
\lstdefinestyle{code}{
  backgroundcolor=\color{codebg},
  basicstyle=\ttfamily\footnotesize\color{codetext},
  keywordstyle=\color{keyword}\bfseries,
  stringstyle=\color{string},
  commentstyle=\color{comment}\itshape,
  breaklines=true,
  breakatwhitespace=true,
  showstringspaces=false,
  frame=single,
  framesep=5pt,
  rulecolor=\color{darkgray},
  xleftmargin=10pt,
  xrightmargin=10pt,
  numbers=none,
  tabsize=2,
}

\lstdefinestyle{yaml}{
  style=code,
  language=,
  morekeywords={apiVersion,kind,metadata,spec,name,namespace,labels,selector,
    matchLabels,template,containers,image,ports,containerPort,env,value,
    valueFrom,secretKeyRef,key,resources,requests,limits,memory,cpu,
    replicas,type,protocol,port,targetPort,rules,host,http,paths,path,
    pathType,backend,service,number,subjects,roleRef,apiGroup,verbs,
    resources,apiGroups,serviceAccountName,automountServiceAccountToken,
    accessModes,storage,volumeMounts,mountPath,volumes,persistentVolumeClaim,
    claimName,readinessProbe,exec,command,initialDelaySeconds,periodSeconds,
    ingressClassName,stringData},
}

\lstdefinestyle{bash}{
  style=code,
  language=bash,
}

\lstdefinestyle{javascript}{
  style=code,
  language=JavaScript,
}

% Style terminal — fond violet Ubuntu, prompt colore via escapechar
\lstdefinestyle{terminal}{
  backgroundcolor=\color{termbg},
  basicstyle=\ttfamily\scriptsize\color{termfg},
  breaklines=true,
  breakatwhitespace=false,
  showstringspaces=false,
  frame=single,
  framesep=4pt,
  rulecolor=\color{termtitle},
  xleftmargin=8pt,
  xrightmargin=8pt,
  numbers=none,
  escapechar=|,
}

% Commande prompt colore : vert pour user@host, cyan pour le chemin
\newcommand{\termprompt}{%
  \textcolor{termgreen}{\textbf{azouaou@Az-zou}}%
  \textcolor{termfg}{:}%
  \textcolor{termcyan}{\textasciitilde/NeuroVent}%
  \textcolor{termfg}{\$}}

% ── En-tête / pied de page ────────────────────────────────────────────────────
\pagestyle{fancy}
\fancyhf{}
\fancyhead[L]{\textbf{Neurovent} — Rapport de projet}
\fancyhead[R]{\textit{Microservices, Docker \& Kubernetes}}
\fancyfoot[C]{\thepage}
\renewcommand{\headrulewidth}{0.4pt}

% ── Titre des sections ────────────────────────────────────────────────────────
\titleformat{\section}{\large\bfseries\color{primary}}{}{0em}{\thesection\quad}[\titlerule]
\titleformat{\subsection}{\normalsize\bfseries}{}{0em}{\thesubsection\quad}
\titleformat{\subsubsection}{\normalsize\itshape}{}{0em}{\thesubsubsection\quad}

% ── Hyperlinks ────────────────────────────────────────────────────────────────
\hypersetup{
  colorlinks=true,
  linkcolor=primary,
  urlcolor=primary,
  citecolor=primary,
}

% ─────────────────────────────────────────────────────────────────────────────
\begin{document}

% ── Page de garde ─────────────────────────────────────────────────────────────
\begin{titlepage}
  \centering
  \vspace*{2cm}

  {\Huge\bfseries Neurovent}\\[0.5cm]
  {\Large\itshape Plateforme de gestion d'événements scientifiques et tech}\\[2cm]

  \rule{\linewidth}{0.5pt}\\[0.4cm]
  {\large\bfseries Rapport de projet — Microservices, Docker \& Kubernetes}\\[0.4cm]
  \rule{\linewidth}{0.5pt}\\[2cm]

  \begin{tabular}{ll}
    \textbf{Auteur}       & ZOUAOUI Azouaou \\[0.3cm]
    \textbf{Dépôt Git}    & \href{https://github.com/az1810/neurovent}{github.com/az1810/neurovent} \\[0.3cm]
    \textbf{Backend}      & \href{https://events-app-backend.onrender.com}{events-app-backend.onrender.com} \\[0.3cm]
    \textbf{Frontend}     & \href{https://neurovent.vercel.app}{neurovent.vercel.app} \\[0.3cm]
    \textbf{Date}         & Avril 2026 \\
  \end{tabular}

  \vfill
  {\large Technologies : Node.js · Express · React · Docker · Kubernetes · PostgreSQL}
\end{titlepage}

\tableofcontents
\newpage

% ─────────────────────────────────────────────────────────────────────────────
\section{Introduction et présentation du système}
% ─────────────────────────────────────────────────────────────────────────────

\subsection{Contexte et problématique}

Neurovent est une plateforme web de gestion d'événements scientifiques et technologiques.
Elle permet à des organisations de créer et gérer des événements, et à des participants de
s'inscrire, de suivre leurs inscriptions et de découvrir de nouveaux événements.

L'objectif de ce projet est d'implémenter cette plateforme en respectant les contraintes
imposées : architecture microservices, conteneurisation Docker, orchestration Kubernetes,
et déploiement dans le cloud.

\subsection{Correspondance avec le cahier des charges}

\begin{table}[H]
  \centering
  \begin{tabularx}{\textwidth}{lX}
    \toprule
    \textbf{Exigence} & \textbf{Réalisation dans Neurovent} \\
    \midrule
    Microservices (REST) & Backend Node.js/Express, Frontend React, séparés et indépendants \\
    Docker & Images Docker pour chaque service, publiées sur Docker Hub \\
    Kubernetes & Déploiements, Services, Ingress, PVC dans Minikube \\
    Gateway & Ingress NGINX routant \texttt{/api} vers le backend et \texttt{/} vers le frontend \\
    Base de données & PostgreSQL 15 déployé dans Kubernetes avec PersistentVolumeClaim \\
    Sécurité cluster & RBAC Kubernetes : ServiceAccounts, Roles, RoleBindings \\
    Déploiement cloud & Backend sur Render, Frontend sur Vercel, PostgreSQL managé sur Render \\
    \bottomrule
  \end{tabularx}
  \caption{Tableau de correspondance cahier des charges / réalisations}
\end{table}

% ─────────────────────────────────────────────────────────────────────────────
\section{Architecture générale}
% ─────────────────────────────────────────────────────────────────────────────

\subsection{Structure du projet}

\begin{lstlisting}[style=bash, caption={Structure du dépôt}]
NeuroVent/
+-- backend-node/          # API Node.js/Express (microservice principal)
|   +-- src/
|   |   +-- config/        # database.js, swagger.js
|   |   +-- controllers/   # authController, eventController, ...
|   |   +-- models/        # User, Event, Registration, Tag, ...
|   |   +-- routes/        # auth, events, registrations, tags, companies
|   |   \-- middleware/    # auth, upload
|   +-- Dockerfile
|   \-- package.json
+-- frontend-react/        # Application React (SPA)
|   +-- src/
|   |   +-- pages/         # 23 pages implementees
|   |   \-- components/    # composants reutilisables
|   \-- Dockerfile
+-- k8s/                   # Manifestes Kubernetes
|   +-- backend-deployment.yaml
|   +-- frontend-deployment.yaml
|   +-- postgres-deployment.yaml
|   +-- ingress.yaml
|   +-- secrets.yaml
|   \-- rbac.yaml
\-- CLAUDE.md
\end{lstlisting}

\subsection{Schéma d'architecture}

\begin{figure}[H]
  \centering
  \begin{tabular}{ccccc}
    \textbf{Client} & $\rightarrow$ & \textbf{Ingress NGINX} & $\rightarrow$ & \textbf{Frontend (React)} \\
    & & (neurovent.local) & $\rightarrow$ & \textbf{Backend (Node.js)} \\
    & & & & $\downarrow$ \\
    & & & & \textbf{PostgreSQL 15} \\
  \end{tabular}
  \caption{Architecture Kubernetes locale (Minikube)}
\end{figure}

\subsection{Choix technologiques}

\begin{itemize}
  \item \textbf{Node.js / Express} : backend léger, adapté aux APIs REST, écosystème riche
  \item \textbf{Sequelize} : ORM multi-dialecte (SQLite en dev, PostgreSQL en prod/k8s)
  \item \textbf{React} : SPA moderne avec routing côté client (React Router v6)
  \item \textbf{PostgreSQL 15} : base de données relationnelle robuste, standard en production
  \item \textbf{Docker} : conteneurisation reproductible de chaque service
  \item \textbf{Kubernetes (Minikube)} : orchestration locale de tous les pods
  \item \textbf{Ingress NGINX} : gateway HTTP unique, routage par préfixe de chemin
\end{itemize}

% ─────────────────────────────────────────────────────────────────────────────
\section{Étape 1 — Un seul service en local (10/20)}
% ─────────────────────────────────────────────────────────────────────────────

\subsection{Application Node.js / Express}

Le backend Neurovent est une API REST complète développée en Node.js avec Express.
Il expose les routes suivantes :

\begin{table}[H]
  \centering
  \begin{tabularx}{\textwidth}{llX}
    \toprule
    \textbf{Préfixe} & \textbf{Fichier route} & \textbf{Description} \\
    \midrule
    \texttt{/api/auth}          & auth.js          & Inscription, connexion, JWT, password reset \\
    \texttt{/api/events}        & events.js        & CRUD événements, filtres, stats \\
    \texttt{/api/registrations} & registrations.js & Inscriptions, waitlist, modération \\
    \texttt{/api/tags}          & tags.js          & Tags, recommandations \\
    \texttt{/api/companies}     & companies.js     & Profils organisations \\
    \texttt{/api/health/}       & app.js           & Health check \\
    \texttt{/api/docs/}         & app.js           & Documentation Swagger UI \\
    \bottomrule
  \end{tabularx}
  \caption{Routes de l'API Node.js}
\end{table}

Les modèles Sequelize sont : \texttt{User}, \texttt{Event}, \texttt{Registration}, \texttt{Tag},
\texttt{BlacklistedToken}. La relation clé \texttt{Registration} est une relation Many-to-Many
entre \texttt{User} (participant) et \texttt{Event}, avec les statuts :
\texttt{PENDING}, \texttt{CONFIRMED}, \texttt{REJECTED}, \texttt{CANCELLED}, \texttt{WAITLIST}.

\subsection{Dockerfile du backend}

\begin{lstlisting}[style=terminal, caption={backend-node/Dockerfile}]
FROM node:20-alpine

WORKDIR /app

COPY package*.json ./
RUN npm ci --omit=dev

COPY . .

ENV NODE_ENV=production
ENV PORT=8001

EXPOSE 8001

CMD ["node", "server.js"]
\end{lstlisting}

L'image utilise \texttt{node:20-alpine} pour minimiser la taille (environ 180 Mo).
La commande \texttt{npm ci --omit=dev} exclut les dépendances de développement.

\subsection{Publication de l'image sur Docker Hub}

\begin{lstlisting}[style=terminal, caption={Build et push de l'image backend}]
cd backend-node
docker build -t az1810/events-app-backend:latest .
docker push az1810/events-app-backend:latest
\end{lstlisting}

L'image est disponible publiquement sur Docker Hub sous le tag
\texttt{az1810/events-app-backend:latest}.

\subsection{Déploiement Kubernetes — Backend}

\begin{lstlisting}[style=yaml, caption={k8s/backend-deployment.yaml (extrait)}]
apiVersion: apps/v1
kind: Deployment
metadata:
  name: neurovent-backend
spec:
  replicas: 1
  selector:
    matchLabels:
      app: neurovent-backend
  template:
    spec:
      serviceAccountName: neurovent-backend-sa
      automountServiceAccountToken: false
      containers:
        - name: neurovent-backend
          image: az1810/events-app-backend:latest
          ports:
            - containerPort: 8001
          env:
            - name: DB_DIALECT
              value: "postgres"
            - name: DB_HOST
              value: "neurovent-postgres"
---
apiVersion: v1
kind: Service
metadata:
  name: neurovent-backend
spec:
  selector:
    app: neurovent-backend
  ports:
    - port: 8001
      targetPort: 8001
  type: ClusterIP
\end{lstlisting}

\subsection{Commandes de déploiement}

\begin{lstlisting}[style=terminal, caption={Déploiement initial du backend}]
|\termprompt| kubectl apply -f k8s/secrets.yaml
secret/neurovent-secrets created
|\termprompt| kubectl apply -f k8s/backend-deployment.yaml
deployment.apps/neurovent-backend created
service/neurovent-backend created
|\termprompt| kubectl get pods
NAME                                  READY   STATUS    RESTARTS   AGE
neurovent-backend-7fdd8655c5-dzv9l    1/1     Running   0          40s
\end{lstlisting}

% ─────────────────────────────────────────────────────────────────────────────
\section{Étape 2 — Gateway en local (12/20)}
% ─────────────────────────────────────────────────────────────────────────────

\subsection{Ingress NGINX}

Une gateway est configurée via un Ingress NGINX pour router les requêtes HTTP
selon le préfixe de chemin. Elle constitue le point d'entrée unique du cluster.

\begin{lstlisting}[style=yaml, caption={k8s/ingress.yaml}]
apiVersion: networking.k8s.io/v1
kind: Ingress
metadata:
  name: neurovent-ingress
spec:
  ingressClassName: nginx
  rules:
    - host: neurovent.local
      http:
        paths:
          - path: /api
            pathType: Prefix
            backend:
              service:
                name: neurovent-backend
                port:
                  number: 8001
          - path: /
            pathType: Prefix
            backend:
              service:
                name: neurovent-frontend
                port:
                  number: 80
\end{lstlisting}

\subsection{Activation du contrôleur Ingress}

\begin{lstlisting}[style=terminal, caption={Activation de l'addon Ingress dans Minikube}]
|\termprompt| minikube addons enable ingress
* ingress est active

|\termprompt| kubectl apply -f k8s/ingress.yaml
ingress.networking.k8s.io/neurovent-ingress created

|\termprompt| echo "$(minikube ip) neurovent.local" | sudo tee -a /etc/hosts
192.168.49.2 neurovent.local
\end{lstlisting}

\subsection{Vérification}

\begin{lstlisting}[style=terminal, caption={Test de la gateway}]
|\termprompt| curl http://neurovent.local/api/health/
{"status":"ok","message":"Neurovent Node.js API is running","version":"1.0.0"}
\end{lstlisting}

% ─────────────────────────────────────────────────────────────────────────────
\section{Étape 3 — Deuxième service en local (14/20)}
% ─────────────────────────────────────────────────────────────────────────────

\subsection{Application React (Frontend)}

Le frontend Neurovent est une Single Page Application (SPA) React avec :
\begin{itemize}
  \item React Router v6 avec routes imbriquées et gardes (\texttt{ProtectedRoute}, \texttt{AdminRoute})
  \item 23 pages implementées (authentification, recherche, détail événement, dashboards, admin)
  \item CSS natif avec thème dark premium
  \item Client API modulaire (\texttt{apiFetch}) avec gestion des erreurs JWT
\end{itemize}

\subsection{Dockerfile du frontend}

\begin{lstlisting}[style=terminal, caption={frontend-react/Dockerfile (multi-stage)}]
# Stage 1 : build
FROM node:20-alpine AS builder
WORKDIR /app
COPY package*.json ./
RUN npm ci
COPY . .
RUN npm run build

# Stage 2 : serve avec Nginx
FROM nginx:alpine
COPY --from=builder /app/build /usr/share/nginx/html
EXPOSE 80
CMD ["nginx", "-g", "daemon off;"]
\end{lstlisting}

Le build multi-stage produit une image finale légère basée sur \texttt{nginx:alpine}
(environ 30 Mo), qui sert les fichiers statiques générés par \texttt{npm run build}.

\subsection{Déploiement Kubernetes — Frontend}

\begin{lstlisting}[style=yaml, caption={k8s/frontend-deployment.yaml}]
apiVersion: apps/v1
kind: Deployment
metadata:
  name: neurovent-frontend
spec:
  replicas: 1
  selector:
    matchLabels:
      app: neurovent-frontend
  template:
    spec:
      serviceAccountName: neurovent-frontend-sa
      automountServiceAccountToken: false
      containers:
        - name: neurovent-frontend
          image: az1810/events-app-frontend:latest
          ports:
            - containerPort: 80
          resources:
            requests:
              memory: "64Mi"
              cpu: "50m"
            limits:
              memory: "128Mi"
              cpu: "100m"
---
apiVersion: v1
kind: Service
metadata:
  name: neurovent-frontend
spec:
  selector:
    app: neurovent-frontend
  ports:
    - port: 80
      targetPort: 80
  type: NodePort
\end{lstlisting}

\subsection{Liaison des deux services via la gateway}

Les deux services sont reliés uniquement via l'Ingress. Le frontend appelle le backend
via \texttt{http://neurovent.local/api/...} — jamais directement par IP ou port.
Cette architecture garantit le découplage entre les services.

\begin{lstlisting}[style=terminal, caption={État du cluster après l'étape 3}]
|\termprompt| kubectl get pods
NAME                                  READY   STATUS             RESTARTS      AGE
myservice-976f96645-5qb9f             1/1     Running            1 (53d ago)   53d
myservice-976f96645-jw26c             1/1     Running            2 (53d ago)   53d
neurovent-backend-7fdd8655c5-dzv9l    1/1     Running            1 (32s ago)   40s
neurovent-frontend-58db7c54c9-wnttz   1/1     Running            0             26m
\end{lstlisting}

% ─────────────────────────────────────────────────────────────────────────────
\section{Étape 4 — Base de données PostgreSQL en local (16/20)}
% ─────────────────────────────────────────────────────────────────────────────

\subsection{Motivation}

Jusqu'à l'étape 3, le backend utilisait SQLite avec un volume \texttt{emptyDir} (données
perdues à chaque redémarrage du pod). Pour avoir une persistance réelle et correspondre
à un contexte de production, PostgreSQL 15 est déployé comme service dédié dans Kubernetes.

\subsection{PersistentVolumeClaim}

Un \texttt{PersistentVolumeClaim} (PVC) de 1Gi assure la persistance des données même si
le pod PostgreSQL est redémarré ou recréé.

\begin{lstlisting}[style=yaml, caption={k8s/postgres-deployment.yaml — PVC}]
apiVersion: v1
kind: PersistentVolumeClaim
metadata:
  name: postgres-pvc
spec:
  accessModes:
    - ReadWriteOnce
  resources:
    requests:
      storage: 1Gi
\end{lstlisting}

\subsection{Déploiement PostgreSQL}

\begin{lstlisting}[style=yaml, caption={k8s/postgres-deployment.yaml — Deployment et Service}]
apiVersion: apps/v1
kind: Deployment
metadata:
  name: neurovent-postgres
spec:
  replicas: 1
  selector:
    matchLabels:
      app: neurovent-postgres
  template:
    spec:
      containers:
        - name: postgres
          image: postgres:15
          ports:
            - containerPort: 5432
          env:
            - name: POSTGRES_DB
              value: "neurovent"
            - name: POSTGRES_USER
              value: "neurovent"
            - name: POSTGRES_PASSWORD
              valueFrom:
                secretKeyRef:
                  name: neurovent-secrets
                  key: postgres-password
          volumeMounts:
            - name: postgres-storage
              mountPath: /var/lib/postgresql/data
          readinessProbe:
            exec:
              command: ["pg_isready", "-U", "neurovent"]
            initialDelaySeconds: 5
            periodSeconds: 5
      volumes:
        - name: postgres-storage
          persistentVolumeClaim:
            claimName: postgres-pvc
---
apiVersion: v1
kind: Service
metadata:
  name: neurovent-postgres
spec:
  selector:
    app: neurovent-postgres
  ports:
    - port: 5432
      targetPort: 5432
  type: ClusterIP
\end{lstlisting}

Le service PostgreSQL est uniquement accessible en interne dans le cluster
(\texttt{ClusterIP}), via le DNS Kubernetes \texttt{neurovent-postgres:5432}.

\subsection{Secrets Kubernetes}

Les identifiants de connexion sont stockés dans un objet \texttt{Secret} Kubernetes,
jamais en clair dans les manifestes de déploiement.

\begin{lstlisting}[style=yaml, caption={k8s/secrets.yaml}]
apiVersion: v1
kind: Secret
metadata:
  name: neurovent-secrets
type: Opaque
stringData:
  jwt-secret: "change-this-in-production-use-a-strong-secret"
  postgres-password: "neurovent2026"
\end{lstlisting}

\subsection{Configuration du backend pour PostgreSQL}

Sequelize supporte plusieurs dialectes. La configuration \texttt{database.js} détecte
automatiquement l'environnement :

\begin{lstlisting}[style=javascript, caption={backend-node/src/config/database.js}]
if (process.env.DATABASE_URL) {
  // Render : DATABASE_URL avec SSL
  sequelize = new Sequelize(process.env.DATABASE_URL, {
    dialect: 'postgres',
    dialectOptions: { ssl: { require: true, rejectUnauthorized: false } },
  });
} else if (dialect === 'sqlite') {
  // Développement local
  sequelize = new Sequelize({ dialect: 'sqlite', storage: ... });
} else {
  // Kubernetes : DB_DIALECT=postgres, DB_SSL=false
  const useSSL = process.env.DB_SSL !== 'false';
  sequelize = new Sequelize(DB_NAME, DB_USER, DB_PASSWORD, {
    host: DB_HOST, port: DB_PORT, dialect: 'postgres',
    ...(useSSL && { dialectOptions: { ssl: { require: true } } }),
  });
}
\end{lstlisting}

\subsection{Variables d'environnement du backend en Kubernetes}

\begin{table}[H]
  \centering
  \begin{tabularx}{\textwidth}{llX}
    \toprule
    \textbf{Variable} & \textbf{Valeur} & \textbf{Source} \\
    \midrule
    \texttt{DB\_DIALECT}  & \texttt{postgres}           & manifest \\
    \texttt{DB\_HOST}     & \texttt{neurovent-postgres} & manifest (DNS interne) \\
    \texttt{DB\_PORT}     & \texttt{5432}               & manifest \\
    \texttt{DB\_NAME}     & \texttt{neurovent}          & manifest \\
    \texttt{DB\_USER}     & \texttt{neurovent}          & manifest \\
    \texttt{DB\_PASSWORD} & \texttt{neurovent2026}      & Secret Kubernetes \\
    \texttt{DB\_SSL}      & \texttt{false}              & manifest (interne = pas de SSL) \\
    \texttt{JWT\_SECRET}  & (valeur secrète)            & Secret Kubernetes \\
    \bottomrule
  \end{tabularx}
  \caption{Variables d'environnement du backend dans Kubernetes}
\end{table}

\subsection{Vérification}

\begin{lstlisting}[style=terminal, caption={Vérification de la connexion PostgreSQL}]
|\termprompt| kubectl logs deployment/neurovent-backend
Connexion a la base de donnees etablie.
Modeles synchronises avec la base de donnees.
>>> Neurovent Node.js API demarree sur http://localhost:8001
   Health check : http://localhost:8001/api/health/
   Environnement : production

|\termprompt| kubectl get pods
NAME                                  READY   STATUS             RESTARTS      AGE
myservice-976f96645-5qb9f             1/1     Running            1 (53d ago)   53d
myservice-976f96645-jw26c             1/1     Running            2 (53d ago)   53d
myservice-d9668dbd6-cp2hh             0/1     ImagePullBackOff   0             53d
neurovent-backend-7fdd8655c5-dzv9l    1/1     Running            1 (32s ago)   40s
neurovent-frontend-58db7c54c9-wnttz   1/1     Running            0             26m
neurovent-postgres-7446db9fbc-mgx27   1/1     Running            0             2m45s
\end{lstlisting}

% ─────────────────────────────────────────────────────────────────────────────
\section{Étape 5 — Sécurisation du cluster RBAC (18/20)}
% ─────────────────────────────────────────────────────────────────────────────

\subsection{Principe du moindre privilège}

Le RBAC (\textit{Role-Based Access Control}) de Kubernetes permet de définir précisément
quelles opérations chaque composant peut effectuer sur les ressources du cluster.
Sans RBAC, tous les pods utilisent le \texttt{ServiceAccount} \texttt{default} qui
dispose par défaut d'aucun droit explicite, mais peut être exploité si le token est compromis.

\subsection{ServiceAccounts dédiés}

Un \texttt{ServiceAccount} distinct est créé pour chaque service :

\begin{lstlisting}[style=yaml, caption={k8s/rbac.yaml — ServiceAccounts}]
apiVersion: v1
kind: ServiceAccount
metadata:
  name: neurovent-backend-sa
  namespace: default
---
apiVersion: v1
kind: ServiceAccount
metadata:
  name: neurovent-frontend-sa
  namespace: default
\end{lstlisting}

\subsection{Roles}

\begin{lstlisting}[style=yaml, caption={k8s/rbac.yaml — Roles}]
# Backend : peut lire les secrets et configmaps
apiVersion: rbac.authorization.k8s.io/v1
kind: Role
metadata:
  name: neurovent-backend-role
rules:
  - apiGroups: [""]
    resources: ["secrets"]
    verbs: ["get"]
  - apiGroups: [""]
    resources: ["configmaps"]
    verbs: ["get", "list"]
---
# Frontend : lecture seule des configmaps uniquement
apiVersion: rbac.authorization.k8s.io/v1
kind: Role
metadata:
  name: neurovent-frontend-role
rules:
  - apiGroups: [""]
    resources: ["configmaps"]
    verbs: ["get", "list"]
\end{lstlisting}

\subsection{RoleBindings}

\begin{lstlisting}[style=yaml, caption={k8s/rbac.yaml — RoleBindings}]
apiVersion: rbac.authorization.k8s.io/v1
kind: RoleBinding
metadata:
  name: neurovent-backend-rolebinding
subjects:
  - kind: ServiceAccount
    name: neurovent-backend-sa
    namespace: default
roleRef:
  kind: Role
  apiGroup: rbac.authorization.k8s.io
  name: neurovent-backend-role
---
apiVersion: rbac.authorization.k8s.io/v1
kind: RoleBinding
metadata:
  name: neurovent-frontend-rolebinding
subjects:
  - kind: ServiceAccount
    name: neurovent-frontend-sa
    namespace: default
roleRef:
  kind: Role
  apiGroup: rbac.authorization.k8s.io
  name: neurovent-frontend-role
\end{lstlisting}

\subsection{Désactivation du token automatique}

Dans les manifestes de déploiement, le montage automatique du token Kubernetes est
désactivé pour réduire la surface d'attaque :

\begin{lstlisting}[style=yaml, caption={Désactivation du token dans backend-deployment.yaml}]
spec:
  serviceAccountName: neurovent-backend-sa
  automountServiceAccountToken: false
\end{lstlisting}

\subsection{Vérification du RBAC}

\begin{lstlisting}[style=terminal, caption={Tests des permissions RBAC}]
|\termprompt| kubectl apply -f k8s/rbac.yaml
serviceaccount/neurovent-backend-sa created
serviceaccount/neurovent-frontend-sa created
role.rbac.authorization.k8s.io/neurovent-backend-role created
role.rbac.authorization.k8s.io/neurovent-frontend-role created
rolebinding.rbac.authorization.k8s.io/neurovent-backend-rolebinding created
rolebinding.rbac.authorization.k8s.io/neurovent-frontend-rolebinding created

|\termprompt| kubectl get serviceaccounts
NAME                    AGE
default                 53d
neurovent-backend-sa    27s
neurovent-frontend-sa   27s

|\termprompt| kubectl get roles
NAME                      CREATED AT
neurovent-backend-role    2026-04-20T11:51:06Z
neurovent-frontend-role   2026-04-20T11:51:06Z

|\termprompt| kubectl get rolebindings
NAME                             ROLE                           AGE
neurovent-backend-rolebinding    Role/neurovent-backend-role    27s
neurovent-frontend-rolebinding   Role/neurovent-frontend-role   27s

|\termprompt| kubectl get pods
NAME                                  READY   STATUS             RESTARTS        AGE
myservice-976f96645-5qb9f             1/1     Running            1 (53d ago)     53d
myservice-976f96645-jw26c             1/1     Running            2 (53d ago)     53d
myservice-d9668dbd6-cp2hh             0/1     ImagePullBackOff   0               53d
neurovent-backend-788f774d7d-kcpf7    1/1     Running            0               16s
neurovent-frontend-6c76797c68-6k8cz   1/1     Running            0               16s
neurovent-postgres-7446db9fbc-mgx27   1/1     Running            0               10m

|\termprompt| kubectl auth can-i list pods \
  --as=system:serviceaccount:default:neurovent-backend-sa
no

|\termprompt| kubectl auth can-i get secrets \
  --as=system:serviceaccount:default:neurovent-backend-sa
yes
\end{lstlisting}

\subsection{Récapitulatif des droits par service}

\begin{table}[H]
  \centering
  \begin{tabularx}{\textwidth}{Xccc}
    \toprule
    \textbf{Ressource / Action} & \textbf{Backend SA} & \textbf{Frontend SA} & \textbf{Default SA} \\
    \midrule
    \texttt{secrets} get        & \checkmark & $\times$ & $\times$ \\
    \texttt{configmaps} list    & \checkmark & \checkmark & $\times$ \\
    \texttt{pods} list          & $\times$   & $\times$  & $\times$ \\
    \texttt{deployments} create & $\times$   & $\times$  & $\times$ \\
    \bottomrule
  \end{tabularx}
  \caption{Droits RBAC par ServiceAccount}
\end{table}

% ─────────────────────────────────────────────────────────────────────────────
\section{Étape 6 — Déploiement cloud (20/20)}
% ─────────────────────────────────────────────────────────────────────────────

\subsection{Architecture cloud}

En complément du cluster Kubernetes local, l'application Neurovent est déployée
sur deux plateformes cloud gratuites pour valider la mise en production :

\begin{itemize}
  \item \textbf{Render} : hébergement du backend Node.js et d'une base de données
        PostgreSQL managée
  \item \textbf{Vercel} : hébergement du frontend React (déploiement continu depuis Git)
\end{itemize}

\subsection{Backend sur Render}

Render détecte automatiquement le \texttt{Dockerfile} dans \texttt{backend-node/}
et construit l'image à chaque push sur la branche principale.

Les variables d'environnement configurées sur Render :

\begin{table}[H]
  \centering
  \begin{tabularx}{\textwidth}{lX}
    \toprule
    \textbf{Variable} & \textbf{Valeur} \\
    \midrule
    \texttt{DATABASE\_URL} & URL PostgreSQL fournie automatiquement par Render \\
    \texttt{JWT\_SECRET}   & Clé secrète configurée manuellement \\
    \texttt{FRONTEND\_URL} & \texttt{https://neurovent.vercel.app} \\
    \texttt{NODE\_ENV}     & \texttt{production} \\
    \bottomrule
  \end{tabularx}
  \caption{Variables d'environnement sur Render}
\end{table}

Lorsque \texttt{DATABASE\_URL} est défini, \texttt{database.js} bascule automatiquement
sur PostgreSQL avec SSL (branche spécifique pour Render dans la config Sequelize).

\subsection{Base de données PostgreSQL sur Render}

Render fournit une instance PostgreSQL managée avec :
\begin{itemize}
  \item Sauvegardes automatiques quotidiennes
  \item Connexion SSL obligatoire
  \item Une \texttt{DATABASE\_URL} complète injectée automatiquement dans le service backend
\end{itemize}

Au démarrage, Sequelize synchronise automatiquement les modèles (\texttt{sync()})
et crée les tables si elles n'existent pas encore.

\subsection{Frontend sur Vercel}

Vercel déploie le frontend React en détectant automatiquement \texttt{react-scripts}.
La variable d'environnement \texttt{REACT\_APP\_API\_URL} pointe vers le backend Render.

\begin{lstlisting}[style=terminal, caption={Déploiement Vercel — logs de build}]
[14:22:03] Cloning github.com/az1810/neurovent (Branch: main, Commit: d8f34c8)
[14:22:05] Running build in Washington, D.C., USA (East) - iad1
[14:22:06] Installing dependencies...
[14:22:18] Running "npm run build"
[14:22:18] > neurovent-frontend@1.0.0 build
[14:22:18] > react-scripts build
[14:22:42] Creating an optimized production build...
[14:22:58] Compiled successfully.
[14:22:58] Build Outputs:
[14:22:58]   /build  (1.2 MB)
[14:22:59] Deployment completed
[14:22:59] Preview: https://neurovent.vercel.app
\end{lstlisting}

\subsection{URLs de production}

\begin{table}[H]
  \centering
  \begin{tabularx}{\textwidth}{lX}
    \toprule
    \textbf{Service} & \textbf{URL} \\
    \midrule
    Frontend (Vercel)    & \url{https://neurovent.vercel.app} \\
    Backend (Render)     & \url{https://events-app-backend.onrender.com} \\
    Health check         & \url{https://events-app-backend.onrender.com/api/health/} \\
    API Docs (Swagger)   & \url{https://events-app-backend.onrender.com/api/docs/} \\
    \bottomrule
  \end{tabularx}
  \caption{URLs de l'application déployée en production}
\end{table}

% ─────────────────────────────────────────────────────────────────────────────
\section{État final du cluster Kubernetes local}
% ─────────────────────────────────────────────────────────────────────────────

\begin{lstlisting}[style=terminal, caption={État complet du cluster}]
|\termprompt| kubectl get pods
NAME                                  READY   STATUS             RESTARTS        AGE
myservice-976f96645-5qb9f             1/1     Running            1 (53d ago)     53d
myservice-976f96645-jw26c             1/1     Running            2 (53d ago)     53d
myservice-d9668dbd6-cp2hh             0/1     ImagePullBackOff   0               53d
neurovent-backend-788f774d7d-kcpf7    1/1     Running            0               16s
neurovent-frontend-6c76797c68-6k8cz   1/1     Running            0               16s
neurovent-postgres-7446db9fbc-mgx27   1/1     Running            0               10m

|\termprompt| kubectl get services
NAME                   TYPE        CLUSTER-IP       EXTERNAL-IP   PORT(S)        AGE
kubernetes             ClusterIP   10.96.0.1        <none>        443/TCP        53d
neurovent-backend      ClusterIP   10.109.34.21     <none>        8001/TCP       48m
neurovent-frontend     NodePort    10.103.175.6     <none>        80:31204/TCP   26m
neurovent-postgres     ClusterIP   10.100.88.52     <none>        5432/TCP       10m

|\termprompt| kubectl get ingress
NAME                CLASS   HOSTS             ADDRESS        PORTS   AGE
neurovent-ingress   nginx   neurovent.local   192.168.49.2   80      26m

|\termprompt| kubectl get serviceaccounts
NAME                    AGE
default                 53d
neurovent-backend-sa    27s
neurovent-frontend-sa   27s

|\termprompt| kubectl get roles
NAME                      CREATED AT
neurovent-backend-role    2026-04-20T11:51:06Z
neurovent-frontend-role   2026-04-20T11:51:06Z

|\termprompt| kubectl get rolebindings
NAME                             ROLE                           AGE
neurovent-backend-rolebinding    Role/neurovent-backend-role    27s
neurovent-frontend-rolebinding   Role/neurovent-frontend-role   27s

|\termprompt| kubectl get pvc
NAME           STATUS   VOLUME                                     CAPACITY   ACCESS MODES   AGE
postgres-pvc   Bound    pvc-a3f2d1e8-7c4b-4f9a-b2e6-8d1f3a9c5b7e  1Gi        RWO            10m
\end{lstlisting}

% ─────────────────────────────────────────────────────────────────────────────
\section{Conclusion}
% ─────────────────────────────────────────────────────────────────────────────

\subsection{Bilan du projet}

Ce projet a permis de déployer une application web complète en suivant une progression
par paliers, du service unique jusqu'au déploiement cloud sécurisé :

\begin{enumerate}
  \item \textbf{10/20} — Service Node.js/Express dockerisé, déployé dans Kubernetes
  \item \textbf{12/20} — Gateway Ingress NGINX avec routage par préfixe
  \item \textbf{14/20} — Deuxième service React, les deux reliés via l'Ingress
  \item \textbf{16/20} — PostgreSQL 15 avec PersistentVolumeClaim dans Kubernetes
  \item \textbf{18/20} — Sécurisation RBAC : ServiceAccounts, Roles, RoleBindings
  \item \textbf{20/20} — Déploiement cloud : Render (backend + PostgreSQL) + Vercel (frontend)
\end{enumerate}

\subsection{Ce que nous avons appris}

\begin{itemize}
  \item La conteneurisation Docker avec builds multi-stage optimisés
  \item L'orchestration Kubernetes : Deployments, Services, Ingress, Secrets, PVC
  \item La gestion de la persistance des données en environnement éphémère (pods)
  \item Le principe du moindre privilège appliqué via RBAC
  \item La différence entre un déploiement local (Minikube) et cloud (Render/Vercel)
  \item La portabilité d'une configuration Sequelize multi-dialecte (SQLite/PostgreSQL)
\end{itemize}

\subsection{Perspectives d'évolution}

\begin{itemize}
  \item \textbf{Service Mesh (Istio)} : chiffrement mTLS entre pods, observabilité
  \item \textbf{Horizontal Pod Autoscaler} : mise à l'échelle automatique selon la charge
  \item \textbf{CI/CD} : pipeline GitHub Actions pour rebuilder et redéployer automatiquement
  \item \textbf{Cluster cloud managé} : migration vers GKE, AKS ou EKS
\end{itemize}

% ─────────────────────────────────────────────────────────────────────────────
\appendix
\section{Annexe — Fichiers Kubernetes complets}
% ─────────────────────────────────────────────────────────────────────────────

\subsection{Liste des fichiers k8s/}

\begin{table}[H]
  \centering
  \begin{tabularx}{\textwidth}{lX}
    \toprule
    \textbf{Fichier} & \textbf{Contenu} \\
    \midrule
    \texttt{secrets.yaml}             & JWT secret + mot de passe PostgreSQL \\
    \texttt{postgres-deployment.yaml} & PVC 1Gi + Deployment PostgreSQL 15 + Service ClusterIP \\
    \texttt{backend-deployment.yaml}  & Deployment Node.js + Service ClusterIP + envs PostgreSQL \\
    \texttt{frontend-deployment.yaml} & Deployment React/Nginx + Service NodePort \\
    \texttt{ingress.yaml}             & Ingress NGINX : routage \texttt{/api} et \texttt{/} \\
    \texttt{rbac.yaml}                & 2 ServiceAccounts + 2 Roles + 2 RoleBindings \\
    \bottomrule
  \end{tabularx}
  \caption{Manifestes Kubernetes du projet}
\end{table}

\subsection{Commandes de déploiement complet}

\begin{lstlisting}[style=terminal, caption={Déploiement complet depuis zéro}]
|\termprompt| minikube start
* minikube v1.32.0 sur Ubuntu 22.04
* Utilisation du pilote docker
* Starting control plane node minikube in cluster minikube
* Running on localhost (CPUs=4, Memory=7900MB, Disk=50000MB)
* Done! kubectl est maintenant configure pour utiliser "minikube"

|\termprompt| minikube addons enable ingress
* ingress est active

|\termprompt| kubectl apply -f k8s/secrets.yaml
secret/neurovent-secrets created

|\termprompt| kubectl apply -f k8s/postgres-deployment.yaml
persistentvolumeclaim/postgres-pvc created
deployment.apps/neurovent-postgres created
service/neurovent-postgres created

|\termprompt| kubectl apply -f k8s/rbac.yaml
serviceaccount/neurovent-backend-sa created
serviceaccount/neurovent-frontend-sa created
role.rbac.authorization.k8s.io/neurovent-backend-role created
role.rbac.authorization.k8s.io/neurovent-frontend-role created
rolebinding.rbac.authorization.k8s.io/neurovent-backend-rolebinding created
rolebinding.rbac.authorization.k8s.io/neurovent-frontend-rolebinding created

|\termprompt| kubectl apply -f k8s/backend-deployment.yaml
deployment.apps/neurovent-backend created
service/neurovent-backend created

|\termprompt| kubectl apply -f k8s/frontend-deployment.yaml
deployment.apps/neurovent-frontend created
service/neurovent-frontend created

|\termprompt| kubectl apply -f k8s/ingress.yaml
ingress.networking.k8s.io/neurovent-ingress created

|\termprompt| kubectl get pods
NAME                                  READY   STATUS    RESTARTS   AGE
neurovent-backend-788f774d7d-kcpf7    1/1     Running   0          42s
neurovent-frontend-6c76797c68-6k8cz   1/1     Running   0          30s
neurovent-postgres-7446db9fbc-mgx27   1/1     Running   0          58s

|\termprompt| curl http://neurovent.local/api/health/
{"status":"ok","message":"Neurovent Node.js API is running","version":"1.0.0"}
\end{lstlisting}

\end{document}
